\documentclass[conference]{IEEEtran}
\IEEEoverridecommandlockouts

\usepackage{cite}
\usepackage{amsmath,amssymb,amsfonts}
\usepackage{graphicx}
\usepackage{textcomp}
\usepackage{xcolor}
\usepackage{booktabs}
\usepackage{multirow}
\usepackage{microtype}

\def\BibTeX{{\rm B\kern-.05em{\sc i\kern-.025em b}\kern-.08em
    T\kern-.1667em\lower.7ex\hbox{E}\kern-.125emX}}

\begin{document}

\title{MinImageScorer: Error-Driven Difficulty Scoring\\
for Guided Copy-Paste Augmentation in\\
Agricultural Object Detection}

\author{
\IEEEauthorblockN{[FILL: Author Name\textsuperscript{1}, Supervisor Name\textsuperscript{1}]}
\IEEEauthorblockA{%
[FILL: Department of XXX, University of XXX, Country]}
}

\maketitle

% ----------------------------------------------------------
\begin{abstract}
We propose \textbf{MinImageScorer}, an error-driven framework that
scores per-object detection difficulty from baseline model ,
combining confidence and localization penalties calibrated via
Pearson correlation optimization. Validated on MinneApple and Global
Wheat Head (GWHD), the score reliably ranks objects and images by
difficulty (Pearson correlation up to $0.82$ with miss rate) and correctly
identifies domain- and size-level failure patterns without supervision. 
This score was then used to guide copy-paste augmentation,
which outperforms random selection on MinneApple
($\Delta\mathrm{AP}{=}{+}0.013$) but both strategies fall below
baseline due to distributional shift. On GWHD, domain shift renders
copy-paste ineffective under either strategy. These results show
that MinImageScorer correctly identifies hard objects, and that
copy-paste can only work with a very limited cased.

(*Does this abstract sound too straight and hard to understand ?
If it is, make it better. do not change the length too much
)
\end{abstract}

\begin{IEEEkeywords}
object detection, copy-paste augmentation, difficulty scoring,
agricultural detection, small objects, domain shift
\end{IEEEkeywords}

% ----------------------------------------------------------
\section{Introduction}

Agricultural object detection benchmarks such as
MinneApple~\cite{hani2020minneapple} and Global Wheat Head
(GWHD)~\cite{david2021global} exhibit distinct structural failure
modes: MinneApple is dominated by sub-resolution small objects, while
GWHD contains severe cross-domain visual variation across geographic
sources. Copy-paste augmentation has demonstrated gains on MS
COCO~\cite{kisantal2019augmentation, ghiasi2021simple}, where objects
are well-distributed and within-distribution. However, the method has
not been thoroughly evaluated on structurally constrained datasets like
agricultural benchmarks, where applying copy-paste blindly risks
corrupting the training distribution rather than improving it. No prior
work provides an empirical framework to diagnose when copy-paste is
beneficial before augmentation is applied. 
(Is this gap addressing good ?
If not do it better, i think we can make it more clear and direct, and also more concise
agriculture dataset have unique failure modes (small objects, domain shift,..) 
I still do not sure about how to connect the gap with the proposed contribution
the contribution is the score that can identify the failure mode.
It was hard to say that using the score to guide copy-paste improve a result
because althought the ap score is increased compared to random but 
it is still below the baseline. and even for gwhd the score-guided was not even better.
We can explain this usign the experiment about the feature space
when just copy and paste has already shift the data too far away from 
the test distribution.
Improve this part for me.
)



We address this gap with three contributions: (1)~\textbf{MinImageScorer}
(MIS), a per-object difficulty score derived from model failure signals
with data-driven weight calibration; (2)~empirical validation showing
the score correctly identifies image-, domain-, and object-size-level
difficulty without access to test labels; (*you can 
think and merge contribution 1,2 together)and (3)~a controlled study
demonstrating that score-guided copy-paste outperforms random selection
on resolution-limited datasets but both strategies fail on
domain-shifted datasets.

% ----------------------------------------------------------
\section{Methodology}

\subsection{Scoring Formulation}

Let $G(I)$ be ground-truth objects in image $I$ and
$\mathcal{P}(g) = \{p : \mathrm{IoU}(p,g) \ge \tau\}$
the matched predictions. Denoting $c_p = \mathrm{conf}(p)$
and $u_p = \mathrm{IoU}(p,g)$, the \emph{object difficulty score} is:

\begin{equation}
S_{\mathrm{obj}}(g) =
\begin{cases}
  \min_{p \in \mathcal{P}(g)}
    \bigl[\alpha(1{-}c_p) + \beta(1{-}u_p)\bigr],
    & \mathcal{P}(g) \neq \varnothing, \\[3pt]
  \alpha + \beta, & \text{otherwise,}
\end{cases}
\label{eq:obj}
\end{equation}

where $\alpha$ and $\beta$ weight confidence and localization penalties.
Missed objects receive the maximum score $\alpha{+}\beta$. The
\emph{image-level score} combines mean object difficulty with
false-positive rate:

\begin{equation}
S_{\mathrm{img}}(I)
  = w_1 \cdot \frac{1}{|G(I)|}\sum_{g\in G(I)} S_{\mathrm{obj}}(g)
  \;+\; w_2 \cdot \mathrm{FPrate}(I).
\label{eq:img}
\end{equation}

Object scoring uses low-confidence predictions ($\tau{=}0.01$) to
capture near-misses; FP rate uses optimal-threshold predictions to
reflect deployment errors. Weights $\alpha, \beta, w_1, w_2$ are
calibrated by maximizing Pearson correlation with per-image miss rate
and FP rate on the training set, requiring no domain-specific
assumptions.

\subsection{Score-Guided Copy-Paste Pipeline}

A baseline detector (YOLOv8s) is trained with all built-in augmentations
disabled. Inference on the training set yields confidence and IoU signals
used to compute $S_{\mathrm{obj}}$ and $S_{\mathrm{img}}$. For
score-guided augmentation, objects are sampled proportionally to
$S_{\mathrm{obj}}$ via exponential weighting; for random, sampling is
uniform. Both conditions paste 2--8 objects per augmented image within
the source image and retrain under identical hyperparameters.

(*Explain more about this part,
I think it too short and not clear enough for reader to image what i do 
with the score guided copy and paste.
)


% ----------------------------------------------------------
\section{Score Validation}
\label{sec:validation}

\subsection{Correlation with Detection Failures}

Table~\ref{tab:corr} reports Pearson correlation between
$S_{\mathrm{img}}$ and per-image miss rate and FP rate on training
images. All cases show strong positive correlation ($r \ge 0.69$),
indicating the score reliably captures per-image difficulty across all
conditions. The strongest relationship (GWHD with traditional
augmentation, $r{=}0.82$) suggests the score accurately identifies hard
images when the baseline model is stable. The inter-seed correlation
remains high across all cases (mean $r \ge 0.80$), confirming that
domain-specific factors drive the score more than training noise.

(*For this experiment i think just keep the the correlation with missrate 
and false positive rate is enough)

\begin{table}[t]
\centering
\caption{Pearson correlation of $S_{\mathrm{img}}$ with per-image miss
rate and FP rate (training set, with traditional augmentation).}
\label{tab:corr}
\resizebox{\columnwidth}{!}{%
\begin{tabular}{@{}lcccc@{}}
\toprule
\textbf{Dataset}
  & $r$(\,score,\,miss\,)
  & $r$(\,score,\,FP\,)
  & Inter-seed $r$ \\
\midrule
MinneApple  & 0.698 & 0.688 & 0.806 \\
GWHD        & 0.816 & 0.818 & 0.900 \\
\bottomrule
\end{tabular}%
}
\end{table}

\subsection{Domain- and Size-Level Validity}

Beyond per-image correlation, the score recovers structural difficulty
without access to test labels. On GWHD, mean $S_{\mathrm{img}}$ exhibits
perfect inverse rank correlation with domain-level AP across all seven
source domains (Spearman $\rho{=}{-}1.00$, $p{<}0.001$). All seven
domains rank identically when ordered by AP and by score: the hardest
domain (\textit{arvalis\_2}, AP\,=\,0.365) scores 0.424, while the
easiest (\textit{rres\_1}, AP\,=\,0.572) scores only 0.200. This
alignment is recovered from training-set model failures alone, without
test labels (Table~\ref{tab:domain}). On MinneApple, small objects
(${\le}32$\,px) average a score of 0.378 vs.\ 0.222 for medium objects,
exactly matching the AP gap (0.132 vs.\ 0.522). Together, these results
confirm that MIS captures image-, domain-, and object-size-level
difficulty from training-set signals without supervision.

(for this part i want to show case the explain the table 2 and one more picture
that show the correlation of the score with the size of object on minneapple
and the score with the domain on gwhd.
State that: the auther of minneapple and gwhd show that the most difficult part 
of this dataset is the small object and the domain shift respectively,
based on the result we just show case, 
we can say that the score can capture this kind of difficulty without any more information.
)

\begin{table}[t]
\centering
\caption{GWHD domain difficulty: mean score vs.\ baseline AP.
  All 7 domains rank identically (Spearman $\rho{=}{-}1.00$,
  $p{<}0.001$, n=7 domains).}
\label{tab:domain}
\resizebox{\columnwidth}{!}{%
\begin{tabular}{@{}lrrrrr@{}}
\toprule
\textbf{Domain} & \textbf{n\,img} & \textbf{AP}
  & \textbf{AP\textsubscript{50}} & \textbf{Mean S\textsubscript{img}}
\\
\midrule
rres\_1    & 363 & 0.572 & 0.962 & 0.200  \\
ethz\_1    & 592 & 0.550 & 0.940 & 0.227  \\
inrae\_1   & 146 & 0.523 & 0.954 & 0.251  \\
arvalis\_3 & 457 & 0.466 & 0.918 & 0.294  \\
usask\_1   & 153 & 0.408 & 0.874 & 0.312  \\
arvalis\_1 & 827 & 0.403 & 0.890 & 0.331  \\
arvalis\_2 & 162 & 0.365 & 0.813 & 0.424  \\
\bottomrule
\end{tabular}%
}
\end{table}

% ----------------------------------------------------------
\section{Augmentation Experiments}

\subsection{Setup}

All models use YOLOv8s initialized from COCO pretrained weights.
Training: MinneApple 200 epochs / patience 10; GWHD 100 epochs /
patience 8; batch 16; input $1024{\times}1024$; lr $10^{-2}$.
All built-in augmentations are disabled to isolate copy-paste effects.
Augmentation ratio 0.3; scale jitter $[0.9,1.1]$ (MA) and
$[0.8,1.2]$ (GWHD).
Do not need to specifiy the hyperparamater too much.
focus more on teh augmentation strategy.
Of we scale the image and rotate it rampdomly
We do that because the researh - cite it show that this is 
a good way of doing copy and paste augmentation.



\subsection{Results}

Table~\ref{tab:ap} shows that on MinneApple, score-guided copy-paste
outperforms random across all sub-metrics ($\Delta\mathrm{AP}
{=}{+}0.013$), confirming the score directs augmentation toward
higher-signal objects. However, both conditions fall below baseline
($-0.024$ for Score CP). The AP$_{75}$ drop ($0.452{\to}0.440$)
indicates degraded localization precision from pasting hard objects
without blending; the AP$_S$ drop ($0.300{\to}0.286$) indicates
score-selected objects are predominantly sub-resolution and cannot be
made learnable by repetition alone.

On GWHD, score-guided and random selection are indistinguishable
($\Delta\mathrm{AP}{=}{-}0.002$) and both fall below baseline
($-0.023$). Score-guided sampling concentrates objects from
domain-minority sources (\textit{arvalis}), amplifying existing domain
imbalance. Backbone feature analysis (YOLOv8s P4, UMAP) confirms that
augmented images are displaced from the test distribution on both
datasets, with greater displacement on GWHD, directly explaining the
consistent degradation.

\begin{table}[t]
\centering
\caption{Three-seed mean COCO AP. \textbf{Bold} = best copy-paste
  condition per dataset.}
\label{tab:ap}
\resizebox{\columnwidth}{!}{%
\begin{tabular}{@{}llrrrrr@{}}
\toprule
\textbf{Dataset} & \textbf{Condition}
  & \textbf{AP}
  & \textbf{AP\textsubscript{50}}
  & \textbf{AP\textsubscript{75}}
  & \textbf{AP\textsubscript{S}}
  & \textbf{AP\textsubscript{M}} \\
\midrule
\multirow{3}{*}{MinneApple}
  & Baseline          & 0.439 & 0.771 & 0.452 & 0.300 & 0.592 \\
  & Random CP         & 0.415 & 0.741 & 0.427 & 0.279 & 0.563 \\
  & \textbf{Score CP} & \textbf{0.428} & \textbf{0.759}
                      & \textbf{0.440} & \textbf{0.286} & \textbf{0.581} \\
\midrule
\multirow{3}{*}{GWHD}
  & Baseline          & 0.507 & 0.922 & 0.491 & 0.136 & 0.501 \\
  & Random CP         & 0.486 & 0.905 & 0.456 & 0.118 & 0.480 \\
  & \textbf{Score CP} & \textbf{0.484} & \textbf{0.905}
                      & \textbf{0.455} & \textbf{0.098} & \textbf{0.477} \\
\bottomrule
\end{tabular}%
}
\end{table}

% ----------------------------------------------------------
\section{Conclusion}

MinImageScorer reliably quantifies detection difficulty from model
failure signals, validated by strong per-image correlations ($r$ up to
$0.82$ with miss rate) and perfect domain-level rank recovery (Spearman
$\rho{=}{-}1.00$ across all 7 GWHD domains without test labels).
Object-size difficulty is correctly ordered on MinneApple. Score-guided
copy-paste outperforms random selection on MinneApple ($+0.013$~AP) but
both strategies degrade below baseline due to distributional shift
between augmented and test images. On GWHD, domain shift makes
copy-paste ineffective regardless of selection quality.

The practical implication is direct: copy-paste is effective when objects
are well-represented in feature space (COCO-like underrepresentation), but
fails when the failure mode is structural (sub-resolution limits on
MinneApple or domain shift on GWHD). MIS provides the diagnostic criterion:
inspect the size and domain distribution of high-score objects before
applying copy-paste augmentation.

Future work should explore two directions: (1)~\textit{Domain-aware
pasting}: select pasted objects only from training domains similar to the
test distribution, reducing the distributional mismatch that causes
degradation here; and (2)~\textit{Dynamic augmentation}: compute object
selection during training rather than fixing scores upfront, allowing the
method to adapt as the model improves. Both approaches aim to keep
augmented images within the test distribution, addressing the core failure
mode identified in this work.

% ----------------------------------------------------------
\begin{thebibliography}{00}

\bibitem{hani2020minneapple}
N.~Hani, P.~Roy, and B.~Bhanu,
``MinneApple: A Benchmark Dataset for Apple Detection and
Segmentation,''
\textit{IEEE Robotics and Automation Letters},
vol.~5, no.~2, pp.~852--858, 2020.

\bibitem{david2021global}
E.~David \textit{et~al.},
``Global Wheat Head Detection (GWHD) Dataset,''
\textit{Plant Phenomics}, vol.~2021, 2021.

\bibitem{kisantal2019augmentation}
M.~Kisantal \textit{et~al.},
``Augmentation for Small Object Detection,''
in \textit{Proc.\ ICACCI}, 2019.

\bibitem{ghiasi2021simple}
G.~Ghiasi \textit{et~al.},
``Simple Copy-Paste is a Strong Data Augmentation Method
for Instance Segmentation,''
in \textit{Proc.\ IEEE/CVF CVPR}, pp.~2918--2928, 2021.

\end{thebibliography}

\end{document}
